% Generated by tools/edn_latex.py from reports/edn.md.
% Do not edit: change reports/edn.md and run `make edn`.
\ednentry{
  id = {EDN-64},
  title = {Pynblint runs from its own locked environment, and a wrapper that reads its JSON report is the gate},
  date = {2026-10-01},
  milestone = {M3: Quality Assurance},
  topic = {Static Analysis, CI/CD},
  participants = {@lukas2510, decided by the repository owner},
  decision = {Pynblint 0.1.6 runs from a dedicated tool environment in \texttt{tools/\allowbreak{}pynblint-\allowbreak{}env/\allowbreak{}}, with its own \texttt{pyproject.\allowbreak{}toml} (\texttt{pynblint=\allowbreak{}=\allowbreak{}0.\allowbreak{}1.\allowbreak{}6}, \texttt{click<8.\allowbreak{}2}, Python 3.12) and its own hash-verified \texttt{uv.\allowbreak{}lock}, invoked as \texttt{uv run -\allowbreak{}-\allowbreak{}project tools/\allowbreak{}pynblint-\allowbreak{}env -\allowbreak{}-\allowbreak{}locked}.
A stdlib-only wrapper, \texttt{tools/\allowbreak{}notebook\_\allowbreak{}lint.\allowbreak{}py}, runs inside that environment, lints each notebook, reads Pynblint's JSON report and is what fails the build, because Pynblint itself exits 0 whatever it finds.
\texttt{make notebook-\allowbreak{}lint} is the one command, and the CI job \texttt{Notebook lint (Pynblint)} runs exactly that target.},
  rationale = {\emph{Alternatives considered:}
\begin{itemize}[leftmargin=*, topsep=2pt]
\item \textbf{Option A (chosen): a locked tool environment, plus a wrapper that is the gate.}
\newline Pros: all 63 packages are pinned and their hashes verified, so the tool cannot change under us; \texttt{-\allowbreak{}-\allowbreak{}locked} fails when the \texttt{pyproject.\allowbreak{}toml} and the lock drift apart; the root \texttt{uv.\allowbreak{}lock} is untouched and the directory is not a workspace member; and the lock is a file Dependabot can read.
The wrapper runs Pynblint with an allowlisted environment in an empty working directory, so no stray variable or \texttt{.\allowbreak{}pynblint} file can reconfigure it, and it fails loudly when Pynblint crashes instead of reading a crash as a clean run.
\newline Cons: two more files, one of them a generated 75 KB lock, a wrapper of its own with tests, and a deviation from the issue's literal ``via \texttt{uvx pynblint}''.
\item \textbf{Option B: \texttt{uvx -\allowbreak{}-\allowbreak{}python 3.\allowbreak{}12 -\allowbreak{}-\allowbreak{}with \textquotesingle{}click=\allowbreak{}=\allowbreak{}8.\allowbreak{}1.\allowbreak{}8\textquotesingle{} pynblint=\allowbreak{}=\allowbreak{}0.\allowbreak{}1.\allowbreak{}6}.}
\newline Pros: one line, and the closest to what the issue asked for.
\newline Cons: about 59 transitive dependencies float, and that is exactly how the tool broke: a click release nobody chose crashed it, so CI can turn red on a pull request that changed nothing near it, and a run cannot be reproduced later.
\item \textbf{Option C: the same, plus \texttt{-\allowbreak{}-\allowbreak{}exclude-\allowbreak{}newer <date>}.}
\newline Pros: the whole resolution is frozen in time, still in one line.
\newline Cons: no hash verification, a magic date in the Makefile, and no lock file for the dependency graph or Dependabot to see.
\item \textbf{Option D: a hashed requirements file passed to \texttt{uvx -\allowbreak{}c}.}
\newline Pros: looks like a lock without a second project.
\newline Cons: \texttt{uvx} silently ignores the hashes; a deliberately tampered hash still ran with exit 0, while \texttt{uv pip install -\allowbreak{}r} on the same file failed with ``Hash mismatch''.
\item \textbf{Option E: a pre-commit hook with \texttt{additional\_\allowbreak{}dependencies}.}
\newline Pros: feedback at commit time, and CI's existing lint job would run it.
\newline Cons: the dependencies float as in option B, and the raw hook reported ``Passed'' on a notebook with findings, so it needs the same wrapper anyway.
\item \textbf{Option F: a dependency group in the root project, kept apart with uv's \texttt{conflicts}.}
\newline Pros: one lock file for everything.
\newline Cons: it ties our lock to an unmaintained tool's \texttt{typer<0.\allowbreak{}13} and \texttt{ipython<9}, and every resolution of the project would carry that weight.
\item \textbf{For the gate: \texttt{jq -\allowbreak{}e} over the JSON in the Makefile, or the wrapper run in the project environment.}
\newline Pros: \texttt{jq} needs no code; the project environment would let the wrapper use typer and loguru like \texttt{tools/\allowbreak{}requirement\_\allowbreak{}matrix.\allowbreak{}py}.
\newline Cons: \texttt{jq} is not on every contributor's machine and prints no reasons; the project environment would make a seconds-long CI job install mlflow, lightgbm and DVC just to lint.
\end{itemize}
\par
\emph{Why:} The decision rests on three measured facts about Pynblint 0.1.6, each reproduced rather than read from its documentation.
First, it does not run as the issue proposed: it pins \texttt{typer<0.\allowbreak{}13} but not click, and typer 0.12 crashes on click 8.2 and newer (released 2025-05-10) with \texttt{TypeError: Secondary flag is not valid for non-\allowbreak{}boolean flag}, on Python 3.12, 3.13 and 3.14 alike.
A tool that broke once through a dependency nobody pinned should not be run in a way that lets it break again, which rules out options B and E and makes the lock worth its two files.
Second, it exits 0 with findings: a notebook with eleven findings still returned 0 in both file and directory mode, so a CI step that only runs it can never fail and something has to read the report.
Third, its configuration leaks in from outside: its settings are read from environment variables with no prefix, so an ordinary \texttt{INCLUDE=\allowbreak{}/\allowbreak{}usr/\allowbreak{}include} crashed a run with a validation error, and a \texttt{.\allowbreak{}pynblint} file in the current directory reconfigures it silently.
The wrapper is the smallest thing that answers all three: it pins nothing itself, it isolates the process, and its exit code is the finding count.
It is stdlib only so that it can run in the tool environment, which keeps the CI job free of the project's own dependencies.},
  ai = {Information seeking, Alternative generation, Alternative assessment, Recommendation, Solution generation},
  response = {Accepted},
  assessment = {AI did the experiments the decision rests on: it reproduced the click crash and pinned it to click rather than to the Python version, measured the exit code with findings, found that \texttt{uvx -\allowbreak{}c} ignores hashes by tampering with one, and timed every isolation option cold and warm.
It also found that the issue's own premise, ``run it isolated via \texttt{uvx pynblint}'', did not work as written.
AI then laid out options A to F with a recommendation, and the owner chose A.
AI also wrote the wrapper and its tests, test-first, and two further AI review passes, one on the code and one on the documents, found defects that were fixed before the pull request was opened.
A mutation run showed that no test noticed when the gate stopped passing the exclusions to Pynblint; an enforced repository rule would have been reported as checked although it never runs on a single notebook; and a report of an unexpected shape exited 1, which reads as findings, instead of 2.
Each now has a test.
An independent AI review of the pull request then found two more: the Pynblint answers the tests replay were recorded once and never checked against the real tool, so a release that stopped reporting findings would have passed every test, and a report that could not be written exited 1, which reads as findings.
A test that runs the real Pynblint from the tool environment and a fix for the exit code followed.
The owner reviews the result in the pull request.},
  aievidence = {Claude Code session on 2026-10-01 implementing issue \#41: a research run on Pynblint 0.1.6 in a scratch directory (rule catalogue, configuration, exit codes, isolation options with timings), then an architecture pass that turned the measurements into options with pros and cons; the owner was given the isolation and gate options with these measurements and chose the locked environment with the wrapper.
Claude Code review session on 2026-10-05 of PR \#67, which ran the gate on deliberately broken copies of the notebook and checked the tests' recorded answers against the real Pynblint.},
  evidence = {\href{https://github.com/mlops-2627q1-mds-upc/MLOPS_Recommenditos/blob/main/tools/pynblint-env/pyproject.toml}{\texttt{tools/\allowbreak{}pynblint-\allowbreak{}env/\allowbreak{}pyproject.\allowbreak{}toml}}, \href{https://github.com/mlops-2627q1-mds-upc/MLOPS_Recommenditos/blob/main/tools/notebook_lint.py}{\texttt{tools/\allowbreak{}notebook\_\allowbreak{}lint.\allowbreak{}py}}, \texttt{tests/\allowbreak{}test\_\allowbreak{}notebook\_\allowbreak{}lint.\allowbreak{}py}, the \texttt{Notebook lint (Pynblint)} job in \href{https://github.com/mlops-2627q1-mds-upc/MLOPS_Recommenditos/blob/main/.github/workflows/ci.yml}{\texttt{.\allowbreak{}github/\allowbreak{}workflows/\allowbreak{}ci.\allowbreak{}yml}}; \hyperref[EDN-65]{EDN-65}; \href{https://github.com/mlops-2627q1-mds-upc/MLOPS_Recommenditos/issues/41}{issue \#41}, \href{https://github.com/mlops-2627q1-mds-upc/MLOPS_Recommenditos/pull/67}{PR \#67}.},
}
